\section{Forschungsdesign}
\label{sec:forschungsdesign}

\subsection{Methodischer Gesamtrahmen und Prozessmodell}
\label{subsec:methodischer_gesamtrahmen}

Die systematische Evaluierung und Auswahl von \textit{Commercial Off-The-Shelf} (COTS)-Software stellt eine bildende, hochgradig interdisziplinäre Problemstellung an der Schnittstelle von Wirtschaftsinformatik, Softwaretechnik und betriebswirtschaftlicher Entscheidungstheorie dar. In modernen Unternehmensarchitekturen nimmt die Bedeutung standardisierter Softwareprodukte kontinuierlich zu, da Organisationen versuchen, Entwicklungszeiten zu verkürzen, Implementierungsrisiken zu minimieren und von kontinuierlichen Herstellerinnovationen zu profitieren. Gleichzeitig führt die hohe Frequenz technologischer Entwicklungen sowie die ausgeprägter Heterogenität veröffentlichter Evaluierungs- und Auswahlansätze dazu, dass eine wissenschaftlich fundierte Bestandsaufnahme eines rigiden, transparenten und reproduzierbaren methodischen Rahmens bedarf. Um sowohl eine vollständige Erfassung des aktuellen internationalen Forschungsstandes als auch eine tiefgehende strukturelle Analyse der existierenden Auswahlmethoden zu gewährleisten, wird in der vorliegenden Studienarbeit ein dreisäuliges hybrides Forschungsdesign konzipiert. Dieses verbindet das etablierte Vorgehensmodell für eigenständige systematische Literaturübersichten nach \cite{2NHFUF5N} mit den prozeduralen Berichts- und Transparenzstandards von PRISMA 2020 \cite{BP42WGEK} sowie der Allgemeinen Morphologischen Analyse (GMA) nach \cite{Z6N6V5L9}.

Die Wahl dieses spezifischen methodischen Gesamtrahmens begründet sich aus den wissenschaftstheoretischen und fachspezifischen Anforderungen der Wirtschaftsinformatik. Als anwendungsorientierte Disziplin verlangt die Wirtschaftsinformatik nicht nur eine rein deskriptive Aggregation bestehender Veröffentlichungen, sondern eine konstruktive Aufbereitung, strukturierte Artefaktevaluation und systematische Rekonstruktion methodischer Lösungsräume. Konventionelle, rein narrative Literaturübersichten neigen dazu, Methoden lediglich chronologisch oder thematisch aufzuzählen, ohne deren strukturelle Kompatibilitäten, mathematischen Prämissen oder situativen Eignungsbedingungen formal gegeneinander abzugrenzen. Die Integration der GMA in den Prozess der systematischen Literaturübersicht (\textit{Systematic Literature Review}, SLR) behebt diesen methodischen Mangel, indem qualitative Methodeneigenschaften systematisch in ein mathematisch-strukturelles Parameterfeld überführt werden.

Das prozedurale Fundament bildet das 8-Schritte-Modell von \cite{2NHFUF5N}, das speziell für die Anforderungen der Informationssystemforschung und Wirtschaftsinformatik konzipiert wurde. Es gliedert den Review-Prozess in die Phasen Identifikation des Zweckes (\textit{Identify Purpose}), Erstellung des Protokolls (\textit{Draft Protocol}), praktisches Screening (\textit{Apply Practical Screen}), Literaturrecherche (\textit{Search Literature}), Datenextraktion (\textit{Extract Data}), Qualitätsbewertung (\textit{Appraise Quality}), Synthese (\textit{Synthesize Studies}) sowie die finale Berichterstattung (\textit{Write Review}). Um eine maximale methodische Rigorosität und Nachvollziehbarkeit sicherzustellen, wird dieses Ablaufmodell operational mit dem PRISMA 2020-Standard \cite{BP42WGEK} synchronisiert. PRISMA liefert hierbei die formalen Richtlinien für das Nachvollziehen von Identifikations-, Screening- und Einschlussentscheidungen und stellt über ein formalisiertes Flow-Diagramm (unter Nutzung der R-basierten Validierung nach \cite{MQ9M8Q7L}) sicher, dass Selektionsverzerrungen (\textit{Selection Bias}) systematisch offenlegt und minimiert werden.

Die operative Schnittstelle zwischen der Literaturübersicht und der Allgemeinen Morphologischen Analyse bildet die Datenextraktions- und Synthesephase. Während die SLR primär der empirischen Materialsammlung, Identifikation, Filterung und qualitativen Aufbereitung der Primärstudien dient, fungiert die GMA als analytisches Instrumentarium der Synthesestufe. Durch die Abbildung der extrahierten Methodeneigenschaften auf ein sechsgliedriges morphologisches Feld gelingt es, den theoretischen Raum aller denkbaren COTS-Auswahlarchitekturen systematisch aufzuspannen und mittels \textit{Cross-Consistency Assessment} (CCA) logische, empirische sowie methodische Inkompatibilitäten innerhalb des Forschungsstands aufzudecken.

Dieses integrierte Vorgehen garantiert, dass die wissenschaftlichen Erkenntnisse dieser Studienarbeit nicht nur auf einer repräsentativen und qualitätsgesicherten Literaturbasis beruhen, sondern auch strukturell verallgemeinerbar und entscheidungstheoretisch fundiert sind. Im Folgenden werden die einzelnen Phasen dieses Forschungsdesigns operational detailliert.

\subsection{Operationalisierung der Literaturrecherche}
\label{subsec:operationalisierung_recherche}

Die Ableitung der operativen Such- und Selektionsprozesse erfolgt streng deduktiv aus den Zielsetzungen der Studienarbeit. Das übergeordnete Forschungsziel besteht in der systematischen Erfassung, kritischen Qualitätsbewertung und strukturlogischen Klassifikation existierender wissenschaftlicher Methoden zur Evaluierung und Auswahl von COTS-Software. Aus diesem Gesamtziel werden drei differenzierte Teilforschungsziele (FZ1 bis FZ3) sowie korrespondierende Forschungsfragen (FF1 bis FF3) abgeleitet, die den roten Faden der Untersuchung definieren:

\begin{itemize}
    \item \textbf{Forschungsziel 1 (FZ1 -- Evolutionäre Entwicklung):} Rekonstruktion der historischen und methodischen Genese formaler COTS-Auswahlansätze von frühen multikriteriellen Grundmodellen bis hin zu modernen hybriden Entscheidungsarchitekturen.
    \begin{itemize}
        \item \textit{Forschungsfrage 1 (FF1):} Wie haben sich wissenschaftliche Methoden zur COTS-Softwareauswahl seit den 2000er Jahren entwickelt und welche mathematisch-theoretischen Grundlagen dominieren die internationale Fachliteratur?
    \end{itemize}
    \item \textbf{Forschungsziel 2 (FZ2 -- Systematischer Vergleich \& Stärken/Grenzen):} Kriteriengeleitete Gegenüberstellung der etablierten Auswahlmethodiken bezüglich ihrer strukturellen Eigenschaften, Komplexität, Eignung für Unsicherheit sowie empirischen Validierung.
    \begin{itemize}
        \item \textit{Forschungsfrage 2 (FF2):} Welche morphologischen Profile weisen bestehende Evaluierungsmethoden auf und wo liegen deren spezifische strukturelle Stärken und praxeologische Limitationen?
    \end{itemize}
    \item \textbf{Forschungsziel 3 (FZ3 -- Identifikation von Forschungslücken):} Systematische Identifikation von blinden Flecken, logischen Inkompatibilitäten und bislang unzureichend erforschten Methodenkonfigurationen im Entscheidungsraum.
    \begin{itemize}
        \item \textit{Forschungsfrage 3 (FF3):} Welche methodischen Forschungslücken und strukturellen Widersprüche existieren im aktuellen Forschungsstand bezüglich der Verknüpfung von Entscheidungsverfahren, Unsicherheitsmodellierung und Praxistauglichkeit?
    \end{itemize}
\end{itemize}

Zur Beantwortung dieser Forschungsfragen wird ein synthetischer Mastersuchstring konstruiert. Dieser basiert auf einer dreigliedrigen Booleschen Verknüpfung, welche die Kernaspekte des Forschungsgegenstands abdeckt: (1) den Software-Objektbereich (\textit{COTS / Standardsoftware}), (2) die prozedurale Evaluierungs- und Auswahltätigkeit (\textit{Evaluation / Selection}) sowie (3) den mathematisch-entscheidungstheoretischen Kern (\textit{Multi-Criteria Decision Making / MCDM}). Durch die systematische Verknüpfung von Synonymen, geschneiderten Platzhaltern (Truncation) und Fachbegriffen ergibt sich der folgende universelle Mastersuchstring:

\begin{quote}
\texttt{("COTS" OR "commercial off-the-shelf" OR "commercial off the shelf" OR "packaged software" OR "standard software") AND ("select*" OR "evaluat*" OR "assess*" OR "choos*" OR "choic*" OR "compar*") AND ("MCDM" OR "MCDA" OR "MADM" OR "multi-criteria" OR "multicriteria" OR "AHP" OR "TOPSIS")}
\end{quote}

Da wissenschaftliche Literaturdatenbanken unterschiedliche Abfragesprachen, Syntaxregeln und Feldindizierungen nutzen, muss der Mastersuchstring spezifisch für die ausgewählten Fachdatenbanken der Wirtschaftsinformatik und Softwaretechnik übersetzt werden. Es werden vier führende elektronische Datenbanken eingebunden: IEEE Xplore, ACM Digital Library, EBSCOhost Business Source Premier sowie Springer Nature eBooks. In Tabelle~\ref{tab:datenbanken_syntax} sind die datenbankspezifischen Adaptionen sowie die durchsuchten Datenfelder detailliert dokumentiert.

\begin{table}[htbp]
\centering
\small
\caption{Datenbankspezifische Übersetzungen und Syntaxadaptionen des Mastersuchstrings}
\label{tab:datenbanken_syntax}
\begin{tabularx}{\textwidth}{p{3.2cm} X p{2.8cm}}
\toprule
\textbf{Datenbank} & \textbf{Spezifische Abfragesyntax / Übersetzung} & \textbf{Durchsuchte Felder} \\
\midrule
\textbf{IEEE Xplore} & \texttt{((("Document Title":COTS OR "Document Title":"commercial off-the-shelf" OR "Abstract":COTS) AND ("Abstract":select* OR "Abstract":evaluat*)) AND ("Full Text \& Metadata":MCDM OR "Full Text \& Metadata":AHP OR "Full Text \& Metadata":TOPSIS))} & Title, Abstract, Index Terms \\
\midrule
\textbf{ACM Digital Library} & \texttt{RecordAbstract:("COTS" OR "commercial off-the-shelf" OR "standard software") AND RecordAbstract:("selection" OR "evaluation") AND RecordAbstract:("MCDM" OR "multi-criteria" OR "AHP")} & Title, Abstract, Keywords \\
\midrule
\textbf{EBSCOhost (Business Source Premier)} & \texttt{TI ( COTS OR "commercial off-the-shelf" OR "packaged software" ) AND AB ( select* OR evaluat* OR assess* ) AND KW ( MCDM OR MCDA OR "multi-criteria" OR AHP OR TOPSIS )} & Title (TI), Abstract (AB), Subject Terms (KW) \\
\midrule
\textbf{Springer Nature eBooks} & \texttt{("COTS" OR "commercial off-the-shelf") AND ("software selection" OR "software evaluation") AND ("multi-criteria" OR "MCDM")} & Content Title, Abstract, Chapter Text \\
\bottomrule
\end{tabularx}
\end{table}

Um aus der Gesamtheit der identifizierten Treffer die primär relevanten Studien objektiv zu selektieren, wird ein dreistufiges \textit{Practical Screening} durchgeführt. Dieses basiert auf der strengen Anwendung vordefinierter Einschlusskriterien (IK) und Ausschlusskriterien (AK). In Tabelle~\ref{tab:einschluss_ausschluss} werden diese Kriterien formal gegeneinander abgegrenzt.

\begin{table}[htbp]
\centering
\small
\caption{Einschluss- (IK) und Ausschlusskriterien (AK) des Practical Screenings}
\label{tab:einschluss_ausschluss}
\begin{tabularx}{\textwidth}{p{0.8cm} X p{0.8cm} X}
\toprule
\textbf{ID} & \textbf{Einschlusskriterium (IK)} & \textbf{ID} & \textbf{Ausschlusskriterium (AK)} \\
\midrule
\textbf{IK1} & Die Studie behandelt explizit die Evaluierung, Auswahl oder den systematischen Vergleich von COTS-Software oder Standardsoftwarekomponenten. & \textbf{AK1} & Die Studie befasst sich ausschließlich mit Hardware-, Netzwerkelement- oder reiner Cloud-Infrastrukturauswahl ohne Anwendungssoftwarebezug. \\
\midrule
\textbf{IK2} & Es wird ein formalisiertes, entscheidungstheoretisches oder mathematisches Modell (z. B. MCDM/MADM, AHP, TOPSIS, Fuzzy-Logik) präsentiert oder evaluiert. & \textbf{AK2} & Der Fokus liegt rein auf Eigenentwicklung (\textit{Make}), Softwarearchitektur-Design oder Maintenance ohne Gegenüberstellung von Drittsoftware. \\
\midrule
\textbf{IK3} & Es handelt sich um eine peer-reviewte wissenschaftliche Veröffentlichung (Journal, Konferenzbeitrag, wissenschaftlicher Buchbeitrag). & \textbf{AK3} & Mangelnde methodische Transparenz; rein kommerzielle Marketing-Whitepaper, Blogbeiträge oder nicht-begutachtete Berichte (Graue Literatur). \\
\midrule
\textbf{IK4} & Der vollständige Text der Studie liegt in englischer oder deutscher Sprache vor. & \textbf{AK4} & Volltext nicht verfügbar oder Veröffentlichungssprache weder Englisch noch Deutsch. \\
\bottomrule
\end{tabularx}
\end{table}

Der praktische Selektionsprozess verläuft schrittweise: Zunächst werden alle identifizierten Datensätze in ein zentrales Referenzmanagementsystem importiert und Duplikate sowohl automatisiert als auch manuell bereinigt. Daraufhin erfolgt in Stufe 1 ein Screening von Titel und Abstract auf Basis der Kriterien IK1--IK4 und AK1--AK4. In Stufe 2 werden die verbleibenden Arbeiten im Volltext eingehend wissenschaftlich geprüft. Studiennebenberichte oder mehrfach publizierte Zwischenergebnisse derselben Autoren werden zu einer Hauptstudie aggregiert, um eine künstliche Verzerrung der Stichprobe zu vermeiden.

\subsection{Qualitätsbewertung und Datenextraktion}
\label{subsec:qualitaetsbewertung_datenextraktion}

Um die wissenschaftliche Belastbarkeit, methodische Rigorosität und praktische Relevanz der in die Synthese eingehenden Primärstudien sicherzustellen, wird ein formalisiertes Qualitätsbewertungsschema (\textit{Quality Appraisal Schema}) angewendet. Gemäß den Empfehlungen von \cite{2NHFUF5N} dient das Quality Appraisal nicht nur der möglichen Exklusion qualitativ unzureichender Arbeiten, sondern primär als gewichtender Faktor bei der anschließenden Interpretation der Ergebnisse, der Beurteilung des Evidenzgrades sowie der Abwägung von Praxiseignung versus theoretischer Komplexität.

Das Bewertungsschema umfasst vier operative Qualitätskriterien (QA1 bis QA4). Jedes Kriterium wird auf einer dreistufigen Ordinalskala von 0 bis 2 Punkten bewertet, wodurch eine maximale Gesamtpunktezahl von 8 Punkten pro Studie erreicht werden kann:

\begin{enumerate}
    \item \textbf{QA1 -- Transparenz \& Formalisierung der Methodik:} 
    \begin{itemize}
        \item \textit{0 Punkte:} Die Evaluierungsmethodik wird lediglich verbal oder oberflächlich skizziert; mathematische Formeln, Algorithmen oder konkrete Prozessschritte fehlen vollständig.
        \item \textit{1 Punkt:} Die Methode ist in ihren Grundzügen nachvollziehbar beschrieben, weist jedoch mathematische Lücken, unvollständig definierte Skalierungen oder vage Aggregationsregeln auf.
        \item \textit{2 Punkte:} Vollständige, mathematisch präzise und nachvollziehbare Formalisierung aller Entscheidungsschritte, Formeln, Gewichtungsverfahren und mathematischen Operationen.
    \end{itemize}
    \item \textbf{QA2 -- Kriterienkatalog-Transparenz \& Strukturierung:} 
    \begin{itemize}
        \item \textit{0 Punkte:} Die zur Softwarebewertung genutzten Evaluierungskriterien werden nicht offengelegt, bleiben rein abstrakt oder sind vollkommen willkürlich gewählt.
        \item \textit{1 Punkt:} Es wird eine exemplarische Liste von Kriterien genannt, ohne dass deren Herleitung, Strukturierung oder taxonomische Verankerung klar wird.
        \item \textit{2 Punkte:} Systematische, theoretisch fundierte Herleitung des Kriterienkatalogs (z. B. unter Rückgriff auf internationale Qualitätsstandards wie ISO/IEC 25010) inklusive transparenter Gewichtungs- und Operationalisierungsvorschriften.
    \end{itemize}
    \item \textbf{QA3 -- Empirische Validierung \& Evidenzgrad:} 
    \begin{itemize}
        \item \textit{0 Punkte:} Die Anwendbarkeit der Methode wird ausschließlich an einem fiktiven, synthetischen Minimalbeispiel (\textit{Toy Example}) demonstriert.
        \item \textit{1 Punkt:} Das Modell wird anhand eines illustrativen Prototyps oder einer vereinfachten akademischen Fallstudie erprobt.
        \item \textit{2 Punkte:} Umfangreiche empirische Validierung in einem realen industriellen Unternehmensumfeld, Durchführung eines kontrollierten Experiments oder empirische Expertenevaluierung in realen Auswahlprojekten.
    \end{itemize}
    \item \textbf{QA4 -- Reflexion von Grenzen \& Anwendungskontext:} 
    \begin{itemize}
        \item \textit{0 Punkte:} Es findet keinerlei Diskussion von Methodenrestriktionen, entscheidungstheoretischen Annahmen oder praktischen Anwendungshürden statt.
        \item \textit{1 Punkt:} Limitationen der Methode werden beiläufig oder oberflächlich angesprochen.
        \item \textit{2 Punkte:} Kritische, tiefgehende Reflexion der Gültigkeitsbereiche, des rechnerischen Aufwands, der Subjektivität von Expertenurteilen sowie der organisatorischen Randbedingungen bei einer Praxiseinführung.
    \end{itemize}
\end{enumerate}

Basierend auf der erreichten Gesamtpunktezahl werden die Primärstudien drei diskreten Evidenzstufen (\textit{Rigor Tiers}) zugeordnet:
\begin{itemize}
    \item \textbf{High Rigor \& Relevance (6--8 Punkte):} Höchste methodische und empirische Qualität. Diese Studien bilden das verlässliche Kernfundament für die Bestimmung valider Parameterprofile in der Morphologischen Analyse.
    \item \textbf{Medium Rigor (3--5 Punkte):} Solide wissenschaftliche Arbeiten mit moderaten Einschränkungen in der empirischen Fundierung oder der Kriterientransparenz. Sie werden in die morphologische Analyse einbezogen, jedoch hinsichtlich ihrer Evidenzstärke differenziert bewertet.
    \item \textbf{Low Rigor / Conceptual (0--2 Punkte):} Arbeiten mit schwerwiegenden Transparenz- oder Formalisierungsdefiziten. Diese Studien werden von der detaillierten morphologischen Synthese ausgeschlossen, um eine Verfälschung des Parameterraums zu verhindern.
\end{itemize}

Für alle eingeschlossenen Studien wird eine umfassende, 18-spaltige Datenextraktionsmatrix konstruiert. Diese Matrix gewährleistet eine feingranulare Dokumentation aller relevanten Merkmale und dient als unmittelbare empirische Schnittstelle zur Morphologischen Analyse. Tabelle~\ref{tab:extraktionsmatrix_struktur} skizziert die systematische Aufteilung der 18 Extraktionsspalten.

\begin{table}[htbp]
\centering
\small
\caption{Struktur und Spaltendefinitionen der 18-spaltigen Datenextraktionsmatrix}
\label{tab:extraktionsmatrix_struktur}
\begin{tabularx}{\textwidth}{p{1.2cm} p{3.8cm} X}
\toprule
\textbf{Spalte} & \textbf{Kategorie} & \textbf{Erfasste Inhaltliche Datenpunkte} \\
\midrule
S1--S4 & Meta-Informationen & Studien-ID, Erstautor, Publikationsjahr, Publikationsorgan (Journal/Konferenz/Buchkapitel). \\
\midrule
S5--S10 & Morphologische Parameter ($P_1$--$P_6$) & Codierung der sechs diskreten GMA-Parameterzustände ($P_1$: Methodentyp, $P_2$: MCDM-Kern, $P_3$: Kriterien-Standard, $P_4$: Unsicherheit, $P_5$: Automation, $P_6$: Validierung). \\
\midrule
S11--S14 & Formal-Methodische Details & Spezifische mathematische Algorithmen (z. B. Saaty-Skala, TOPSIS-Distanzmaß), Konsistenzprüfverfahren (z. B. Consistency Ratio CR), Aggregationsoperatoren, Gewichtungsmethoden. \\
\midrule
S15--S17 & Empirische \& Praktische Aspekte & Softwareunterstützung/Tooling (z. B. Excel, MATLAB, dediziertes DSS), Anwendungsdomäne (z. B. ERP, CRM, DBMS, IoT), Rechnerische Komplexität \& Skalierbarkeit. \\
\midrule
S18 & Qualitäts-Score \& Evidenz & Erreichte Punkte in QA1--QA4 (0--8 Punkte) sowie finale Einstufung in Rigor-Tier (High/Medium/Low). \\
\bottomrule
\end{tabularx}
\end{table}

Die Befüllung der Matrix erfolgt nach dem Vier-Augen-Prinzip, wobei Abweichungen in der Codierung durch gegenseitige Diskurse und Referenzierung des Ursprungstextes konsolidiert werden.

\subsection{Methodik der Morphologischen Analyse}
\label{subsec:methodik_morphologische_analyse}

Die Allgemeine Morphologische Analyse (GMA), ursprungsphilosophisch entwickelt durch den Schweizer Astrophysiker Fritz Zwicky und für moderne sozial- und ingenieurwissenschaftliche Systemkontexte von \cite{Z6N6V5L9} weiterentwickelt, stellt ein nicht-quantitatives Modellierungsinstrumentarium zur Erforschung komplexer, mehrdimensionaler Problemräume dar. In der Wirtschaftsinformatik eignet sich die GMA in hervorragender Weise zur Typologisierung von Systemen, Architekturmodellen und Entscheidungsverfahren, da sie es erlaubt, vielschichtige qualitative Zusammenhänge frei von unzulässigen Vereinfachungen, Reduktionismen oder erzwungenen Quantifizierungen formallogisch zu strukturieren.

Die mathematische Grundlage der GMA besteht in der Aufspannung eines diskreten $k$-dimensionalen Parameterraumes (des sogenannten \textit{Morphologischen Feldes}). Sei $P = \{P_1, P_2, \dots, P_k\}$ die Menge von $k$ voneinander unabhängigen, aber strukturell wechselwirkenden Parametern, die den Untersuchungsgegenstand vollständig beschreiben. Jedem Parameter $P_i$ wird eine endliche Menge von $m_i$ disjunkten Zustandswerten $V_i = \{v_{i,1}, v_{i,2}, \dots, v_{i,m_i}\}$ zugeordnet. 

Für den spezifischen Forschungsgegenstand der COTS-Softwareauswahl wird auf Basis der entscheidungstheoretischen Grundlagen und der systematischen Sichtung des Forschungsstandes ein 6-dimensionales morphologisches Feld ($P_1$ bis $P_6$) definiert:

\begin{enumerate}
    \item \textbf{Parameter 1 ($P_1$) -- Methodentypus / Problemstrukturierung:} Beschreibt die grundlegende entscheidungstheoretische Architektur und Denkweise des Ansatzes.
    \begin{itemize}
        \item $v_{1,1}$: Monokriterielle / Nutzwertanalytische Grundmodelle (NWA / Scoring-Modelle)
        \item $v_{1,2}$: Hierarchisch-dekomponierende Strukturmodelle
        \item $v_{1,3}$: Outranking- und Dominanzstrukturmodelle
        \item $v_{1,4}$: Mathematische Zielprogrammierung \& Optimierungsmodelle
    \end{itemize}

    \item \textbf{Parameter 2 ($P_2$) -- MCDM/MADM-Kernverfahren:} Spezifiziert den primär eingesetzten Algorithmus zur Präferenzaggregation und Alternativenreihung.
    \begin{itemize}
        \item $v_{2,1}$: Analytic Hierarchy / Network Process (AHP / ANP)
        \item $v_{2,2}$: Distanzbasierte Verfahren (TOPSIS / VIKOR)
        \item $v_{2,3}$: Paarvergleichs- und Outranking-Verfahren (PROMETHEE / ELECTRE)
        \item $v_{2,4}$: Hybride / Multi-Methoden-Kombinationen
    \end{itemize}

    \item \textbf{Parameter 3 ($P_3$) -- Kriterienkatalog-Standardisierung:} Erfasst den Reifegrad, die Normierung und die Taxonomie der verwendeten Evaluierungskriterien.
    \begin{itemize}
        \item $v_{3,1}$: Ad-hoc / rein projektspezifische Kriterienentwicklung
        \item $v_{3,2}$: Domänenspezifische / branchenbezogene Kriterientaxonomien
        \item $v_{3,3}$: Standardisierte / international normierte Taxonomien (z. B. ISO/IEC 25010)
    \end{itemize}

    \item \textbf{Parameter 4 ($P_4$) -- Unsicherheits- \& Präferenzmodellierung:} Spiegelt den mathematischen Umgang mit vagem, unvollständigem oder subjektivem Expertenwissen wider.
    \begin{itemize}
        \item $v_{4,1}$: Deterministische Modelle (exakte Punktbewertungen)
        \item $v_{4,2}$: Unscharfe Mengen / Fuzzy-Logik (Fuzzy-AHP, Fuzzy-TOPSIS, Intuitionistic Fuzzy)
        \item $v_{4,3}$: Stochastische Modelle \& Wahrscheinlichkeitsnetzwerke (Bayes'sche Netze / Monte-Carlo-Simulationen)
    \end{itemize}

    \item \textbf{Parameter 5 ($P_5$) -- Decision Support Automation Level:} Bezeichnet den Grad der steuerungstechnischen und softwareseitigen Automatisierung des Auswahlprozesses.
    \begin{itemize}
        \item $v_{5,1}$: Manueller Prozess / tabellenkalkulationsbasierte Umsetzung (z. B. Excel)
        \item $v_{5,2}$: Teilautomatisierte Skript- und Expertensystemlösungen
        \item $v_{5,3}$: Vollautomatisierte, plattformintegrierte Decision-Support-Systeme (DSS)
    \end{itemize}

    \item \textbf{Parameter 6 ($P_6$) -- Empirische Validierungsform \& Anwendungsbereich:} Beschreibt das Ausmaß und den Evidenzcharakter der praktischen Erprobung.
    \begin{itemize}
        \item $v_{6,1}$: Pure Konzeption / Demonstratives Minimalbeispiel (\textit{Toy Example})
        \item $v_{6,2}$: Einzelne industrielle Fallstudie / Pilotprojekt
        \item $v_{6,3}$: Empirische Multi-Case-Validierung / Feldstudie
    \end{itemize}
\end{enumerate}

Der theoretisch aufgespannte Gesamtraum $N$ aller denkbaren Methodenkonfigurationen berechnet sich aus dem kartesischen Produkt der Anzahl der Parameterzustände $|v_i|$ gemäß Formel~(\ref{eq:konfigurationsraum}):

\begin{equation}
\label{eq:konfigurationsraum}
N = \prod_{i=1}^6 |v_i| = |v_1| \times |v_2| \times |v_3| \times |v_4| \times |v_5| \times |v_6| = 4 \times 4 \times 3 \times 3 \times 3 \times 3 = 1.296
\end{equation}

Der Gesamtraum umfasst somit exakt 1.296 theoretisch mögliche Methodenprofile. Ein zentraler methodischer Schritt der GMA nach \cite{Z6N6V5L9} besteht nun in der systematischen Reduktion dieses immensen Problemraums auf den Bereich der tatsächlich logisch konsistenten, praxeologisch realisierbaren und empirisch nachgewiesenen Lösungsräume. Dies erfolgt über das sogenannte \textit{Cross-Consistency Assessment} (CCA).

Beim CCA wird das sechsgliedrige morphologische Feld in eine symmetrische Cross-Consistency-Matrix (CCM) überführt. Innerhalb dieser Matrix wird jedes mögliche Paar von Parameterzuständen ($v_{i,a}, v_{j,b}$) mit $i \neq j$ systematisch auf seine Verträglichkeit hin überprüft. Das CCA unterscheidet drei primäre Typen von Inkompatibilitäten beziehungsweise Restriktionen:

\begin{enumerate}
    \item \textbf{Logische Widersprüche (Pure Logical Contradictions):} Verknüpfungen von Zuständen, die sich aus rein formal-logischen oder mathematischen Gründen ausschließen. Beispielsweise ist die Verknüpfung von deterministischer Präferenzmodellierung ($v_{4,1}$) mit rein stochastischen Bayes-Netzwerken ($v_{2,4}$ bzw. $v_{4,3}$) logisch widersprüchlich.
    \item \textbf{Empirische Inkongruenzen (Empirical Mismatches):} Kombinationen, die theoretisch-abstrakt denkbar wären, in der Praxis jedoch aufgrund struktureller Gegensätze scheitern. So steht beispielsweise die Forderung nach vollautomatisierten DSS-Systemen ($v_{5,3}$) im empirischen Widerspruch zu rein ad-hoc generierten, unstrukturierten Kriterienkatalogen ($v_{3,1}$), da letztere keine maschinenlesbare Semantik bieten.
    \item \textbf{Methodische Redundanzen / Normative Constraints:} Kombinationen, die rechnerisch hochgradig suboptimal oder ineffizient sind. Hierzu zählt etwa die Verknüpfung extrem komplexer Fuzzy-Outranking-Verfahren ($v_{2,3}, v_{4,2}$) mit rein manueller Tabellenkalkulation ($v_{5,1}$) bei großen Kriterienmengen, da hier die mathematische Operatorkomplexität die menschliche Verarbeitungskapazität überschreitet.
\end{enumerate}

Durch den systematischen Vollzug des CCA werden sämtliche inkompatible Paare aus der Cross-Consistency-Matrix eliminiert. Die verbleibende Menge konsistenter Konfigurationen repräsentiert das finale Lösungsfeld. Die Gegenüberstellung der im SLR empirisch vorgefundenen Studienmuster mit dem theoretisch reduzierten konsistenten Lösungsraum erlaubt es schließlich, dominante Methodencluster zu identifizieren sowie weiße Flecken -- sprich bislang unbesetzte, aber theoretisch konsistente und vielversprechende Methodenkonfigurationen -- aufzudecken.
